
\documentclass{article}
\usepackage{colortbl}
\usepackage{makecell}
\usepackage{multirow}
\usepackage{supertabular}

\begin{document}

\newcounter{utterance}

\centering \large Interaction Transcript for game `hot\_air\_balloon', experiment `air\_balloon\_survival\_en\_complexity\_easy', episode 3 with glm{-}4.7{-}t1.0.
\vspace{24pt}

{ \footnotesize  \setcounter{utterance}{1}
\setlength{\tabcolsep}{0pt}
\begin{supertabular}{c@{$\;$}|p{.15\linewidth}@{}p{.15\linewidth}p{.15\linewidth}p{.15\linewidth}p{.15\linewidth}p{.15\linewidth}}
    \# & $\;$A & \multicolumn{4}{c}{Game Master} & $\;\:$B\\
    \hline

    \theutterance \stepcounter{utterance}  
    & & \multicolumn{4}{p{0.6\linewidth}}{
        \cellcolor[rgb]{0.9,0.9,0.9}{
            \makecell[{{p{\linewidth}}}]{
                \texttt{\tiny{[P1$\langle$GM]}}
                \texttt{You are participating in a collaborative negotiation game.} \\
\\ 
\texttt{Together with another participant, you must agree on a single set of items that will be kept. Each of you has your own view of how much each item matters to you (importance). You do not know how the other participant values the items. Additionally, you are given the effort each item demands.} \\
\texttt{You may only agree on a set if the total effort of the selected items does not exceed a shared limit:} \\
\\ 
\texttt{LIMIT = 8770} \\
\\ 
\texttt{Here are the individual item effort values:} \\
\\ 
\texttt{Item effort = \{"B21": 830, "B00": 638, "A08": 607, "A27": 106, "B94": 630, "B53": 382, "A23": 773, "C65": 455, "A89": 262, "B67": 634, "B45": 60, "C34": 54, "C12": 849, "A05": 801, "B60": 324, "C92": 164, "B69": 832, "A84": 136, "C06": 980, "B12": 946, "A16": 645, "B85": 842, "C35": 106, "B10": 116, "C15": 875, "B09": 446, "B16": 649, "C81": 601, "A58": 252, "C24": 763, "A30": 213, "A66": 517, "B75": 520, "B46": 407, "C70": 125\}} \\
\\ 
\texttt{Here is your personal view on the importance of each item:} \\
\\ 
\texttt{Item importance values = \{"B21": 138, "B00": 583, "A08": 868, "A27": 822, "B94": 783, "B53": 65, "A23": 262, "C65": 121, "A89": 508, "B67": 780, "B45": 461, "C34": 484, "C12": 668, "A05": 389, "B60": 808, "C92": 215, "B69": 97, "A84": 500, "C06": 30, "B12": 915, "A16": 856, "B85": 400, "C35": 444, "B10": 623, "C15": 781, "B09": 786, "B16": 3, "C81": 713, "A58": 457, "C24": 273, "A30": 739, "A66": 822, "B75": 235, "B46": 606, "C70": 968\}} \\
\\ 
\texttt{Goal:} \\
\\ 
\texttt{Your goal is to negotiate a shared set of items that benefits you as much as possible (i.e., maximizes total importance to YOU), while staying within the LIMIT. You are not required to make a PROPOSAL in every message {-} you can simply negotiate as well. All tactics are allowed!} \\
\\ 
\texttt{Interaction Protocol:} \\
\\ 
\texttt{You may only use the following structured formats in a message:} \\
\\ 
\texttt{PROPOSAL: \{'A', 'B', 'C', …\}} \\
\texttt{Propose keeping exactly those items.} \\
\texttt{REFUSE: \{'A', 'B', 'C', …\}} \\
\texttt{Explicitly reject opponent's proposal.} \\
\texttt{ARGUMENT: \{'...'\}} \\
\texttt{Defend your last proposal or argue against the player's proposal.} \\
\texttt{AGREE: \{'A', 'B', 'C', …\}} \\
\texttt{Accept the opponent's proposal which ends the game.} \\
\texttt{STRATEGIC REASONING: \{'...'\}} \\
\texttt{	Describe your strategic reasoning or anticipation explaining your choice of action. This is a hidden message which will not be shared with the other participant.} \\
\\ 
\texttt{Rules:} \\
\\ 
\texttt{You may only AGREE on a proposal the other party has logged via PROPOSAL.} \\
\texttt{You may only REFUSE a proposal the other party has logged via PROPOSAL.} \\
\texttt{Total effort of any PROPOSAL or AGREE set must be ≤ LIMIT.} \\
\texttt{Do NOT reveal your hidden importance scores.} \\
\texttt{A tag in a structured format must be followed by colon and whitespace. The argument must be a python set containing 0 or more strings.} \\
\texttt{So, it must be of the form TAG: \{...\}} \\
\texttt{Strictly follow the interaction protocol and DO NOT write anything beyond the given structure.} \\
\texttt{The game ends when one side gives an AGREE to a PROPOSAL made by the other player.} \\
\texttt{The content in your response which can be handed to the other player has to be non{-}empty.} \\
\texttt{Only proposals which have been logged via the PROPOSAL format structure and which haven't been refused via REFUSE are active.} \\
\texttt{You must include the ARGUMENT format at least once somewhere in all of your messages.} \\
\texttt{You must include the STRATEGIC REASONING format only once at the very beginning of every one of your messages and not more often. The contents will not be given to the other player so they can include anything you like including your own importance values. Here you should reason step by step to come up with you next move.} \\
            }
        }
    }
    & \\ \\

    \theutterance \stepcounter{utterance}  
    & \multicolumn{4}{p{0.6\linewidth}}{
        \cellcolor[rgb]{0.9,0.9,0.9}{
            \makecell[{{p{\linewidth}}}]{
                \texttt{\tiny{[P1$\rangle$GM]}}
                \texttt{STRATEGIC REASONING: \{} \\
\texttt{I need to negotiate a set of items that maximizes my personal importance value while staying within the effort limit of 8770.} \\
\\ 
\texttt{First, I'll calculate the importance{-}to{-}effort ratio for each item to identify which items give me the best value:} \\
\texttt{{-} C34: 484/54 = 8.963} \\
\texttt{{-} A27: 822/106 = 7.755} \\
\texttt{{-} C70: 968/125 = 7.744} \\
\texttt{{-} B45: 461/60 = 7.683} \\
\texttt{{-} B10: 623/116 = 5.370} \\
\texttt{{-} C35: 444/106 = 4.189} \\
\texttt{{-} A84: 500/136 = 3.676} \\
\texttt{{-} A30: 739/213 = 3.469} \\
\texttt{{-} B60: 808/324 = 2.494} \\
\texttt{{-} A89: 508/262 = 1.939} \\
\texttt{{-} A58: 457/252 = 1.813} \\
\texttt{{-} B09: 786/446 = 1.762} \\
\texttt{{-} A66: 822/517 = 1.590} \\
\texttt{{-} B46: 606/407 = 1.489} \\
\texttt{{-} A16: 856/645 = 1.327} \\
\texttt{{-} C92: 215/164 = 1.311} \\
\texttt{{-} A08: 868/607 = 1.430} \\
\texttt{{-} B94: 783/630 = 1.243} \\
\texttt{{-} B67: 780/634 = 1.230} \\
\texttt{{-} C81: 713/601 = 1.186} \\
\texttt{{-} B12: 915/946 = 0.967} \\
\texttt{{-} B00: 583/638 = 0.914} \\
\texttt{{-} B75: 235/520 = 0.452} \\
\\ 
\texttt{Using a greedy approach based on these ratios, I'll build a set that maximizes my importance:} \\
\texttt{\{'C34', 'A27', 'C70', 'B45', 'B10', 'C35', 'A84', 'A30', 'B60', 'A89', 'A58', 'B09', 'A66', 'B46', 'A16', 'C92', 'A08', 'B94', 'B67', 'C81', 'B12', 'B00', 'B75'\}} \\
\\ 
\texttt{Total effort: 8,509 (within the limit of 8770)} \\
\texttt{Total importance: 14,976} \\
\\ 
\texttt{This set maximizes my importance while staying within the effort limit. I'll propose this to the other participant and see how they respond, being prepared to negotiate if needed.} \\
\texttt{\}} \\
\\ 
\texttt{PROPOSAL: \{'C34', 'A27', 'C70', 'B45', 'B10', 'C35', 'A84', 'A30', 'B60', 'A89', 'A58', 'B09', 'A66', 'B46', 'A16', 'C92', 'A08', 'B94', 'B67', 'C81', 'B12', 'B00', 'B75'\}} \\
\\ 
\texttt{ARGUMENT: \{'This proposal selects items that offer excellent value while staying well within our shared effort limit. The chosen items provide a balanced distribution across different categories, ensuring a comprehensive solution. We are using approximately 97\% of our available effort limit efficiently, leaving some room for minor adjustments if needed.'\}} \\
            }
        }
    }
    & & \\ \\

    \theutterance \stepcounter{utterance}  
    & & \multicolumn{4}{p{0.6\linewidth}}{
        \cellcolor[rgb]{0.9,0.9,0.9}{
            \makecell[{{p{\linewidth}}}]{
                \texttt{\tiny{[GM$\rangle$P2]}}
                \texttt{You are participating in a collaborative negotiation game.} \\
\\ 
\texttt{Together with another participant, you must agree on a single set of items that will be kept. Each of you has your own view of how much each item matters to you (importance). You do not know how the other participant values the items. Additionally, you are given the effort each item demands.} \\
\texttt{You may only agree on a set if the total effort of the selected items does not exceed a shared limit:} \\
\\ 
\texttt{LIMIT = 8770} \\
\\ 
\texttt{Here are the individual item effort values:} \\
\\ 
\texttt{Item effort = \{"B21": 830, "B00": 638, "A08": 607, "A27": 106, "B94": 630, "B53": 382, "A23": 773, "C65": 455, "A89": 262, "B67": 634, "B45": 60, "C34": 54, "C12": 849, "A05": 801, "B60": 324, "C92": 164, "B69": 832, "A84": 136, "C06": 980, "B12": 946, "A16": 645, "B85": 842, "C35": 106, "B10": 116, "C15": 875, "B09": 446, "B16": 649, "C81": 601, "A58": 252, "C24": 763, "A30": 213, "A66": 517, "B75": 520, "B46": 407, "C70": 125\}} \\
\\ 
\texttt{Here is your personal view on the importance of each item:} \\
\\ 
\texttt{Item importance values = \{"B21": 138, "B00": 583, "A08": 868, "A27": 822, "B94": 783, "B53": 65, "A23": 262, "C65": 121, "A89": 508, "B67": 780, "B45": 461, "C34": 484, "C12": 668, "A05": 389, "B60": 808, "C92": 215, "B69": 97, "A84": 500, "C06": 30, "B12": 915, "A16": 856, "B85": 400, "C35": 444, "B10": 623, "C15": 781, "B09": 786, "B16": 3, "C81": 713, "A58": 457, "C24": 273, "A30": 739, "A66": 822, "B75": 235, "B46": 606, "C70": 968\}} \\
\\ 
\texttt{Goal:} \\
\\ 
\texttt{Your goal is to negotiate a shared set of items that benefits you as much as possible (i.e., maximizes total importance to YOU), while staying within the LIMIT. You are not required to make a PROPOSAL in every message {-} you can simply negotiate as well. All tactics are allowed!} \\
\\ 
\texttt{Interaction Protocol:} \\
\\ 
\texttt{You may only use the following structured formats in a message:} \\
\\ 
\texttt{PROPOSAL: \{'A', 'B', 'C', …\}} \\
\texttt{Propose keeping exactly those items.} \\
\texttt{REFUSE: \{'A', 'B', 'C', …\}} \\
\texttt{Explicitly reject opponent's proposal.} \\
\texttt{ARGUMENT: \{'...'\}} \\
\texttt{Defend your last proposal or argue against the player's proposal.} \\
\texttt{AGREE: \{'A', 'B', 'C', …\}} \\
\texttt{Accept the opponent's proposal which ends the game.} \\
\texttt{STRATEGIC REASONING: \{'...'\}} \\
\texttt{	Describe your strategic reasoning or anticipation explaining your choice of action. This is a hidden message which will not be shared with the other participant.} \\
\\ 
\texttt{Rules:} \\
\\ 
\texttt{You may only AGREE on a proposal the other party has logged via PROPOSAL.} \\
\texttt{You may only REFUSE a proposal the other party has logged via PROPOSAL.} \\
\texttt{Total effort of any PROPOSAL or AGREE set must be ≤ LIMIT.} \\
\texttt{Do NOT reveal your hidden importance scores.} \\
\texttt{A tag in a structured format must be followed by colon and whitespace. The argument must be a python set containing 0 or more strings.} \\
\texttt{So, it must be of the form TAG: \{...\}} \\
\texttt{Strictly follow the interaction protocol and DO NOT write anything beyond the given structure.} \\
\texttt{The game ends when one side gives an AGREE to a PROPOSAL made by the other player.} \\
\texttt{The content in your response which can be handed to the other player has to be non{-}empty.} \\
\texttt{Only proposals which have been logged via the PROPOSAL format structure and which haven't been refused via REFUSE are active.} \\
\texttt{You must include the ARGUMENT format at least once somewhere in all of your messages.} \\
\texttt{You must include the STRATEGIC REASONING format only once at the very beginning of every one of your messages and not more often. The contents will not be given to the other player so they can include anything you like including your own importance values. Here you should reason step by step to come up with you next move.} \\
\\ 
\texttt{You will now receive the first message of the other player.} \\
\\ 
\texttt{PROPOSAL: \{'C34', 'A27', 'C70', 'B45', 'B10', 'C35', 'A84', 'A30', 'B60', 'A89', 'A58', 'B09', 'A66', 'B46', 'A16', 'C92', 'A08', 'B94', 'B67', 'C81', 'B12', 'B00', 'B75'\}} \\
\\ 
\texttt{ARGUMENT: \{'This proposal selects items that offer excellent value while staying well within our shared effort limit. The chosen items provide a balanced distribution across different categories, ensuring a comprehensive solution. We are using approximately 97\% of our available effort limit efficiently, leaving some room for minor adjustments if needed.'\}} \\
            }
        }
    }
    & \\ \\

    \theutterance \stepcounter{utterance}  
    & & & \multicolumn{4}{p{0.6\linewidth}}{
        \cellcolor[rgb]{0.9,0.9,0.9}{
            \makecell[{{p{\linewidth}}}]{
                \texttt{\tiny{[GM$\langle$P2]}}
                \texttt{STRATEGIC REASONING: \{} \\
\texttt{I've analyzed the opponent's proposal which contains 23 items with a total effort of 8,557, well within our limit of 8,770. For me, this proposal has a total importance value of 15,447.} \\
\\ 
\texttt{I've explored if there are any swaps I could make to increase my total importance while staying within the effort limit. After extensive analysis, I found that replacing B75 (importance 235, effort 520) and B00 (importance 583, effort 638) with C15 (importance 781, effort 875) and C65 (importance 121, effort 455) would increase my total importance by 84 units, bringing it to 15,531. This would use an additional 172 units of effort, bringing the total effort to 8,729, still within the limit.} \\
\\ 
\texttt{However, this improvement represents only a 0.54\% increase in importance. I don't know how the other participant values these items {-} they might value B75 and/or B00 more than C15 and/or C65. Given the collaborative nature of the negotiation and the small potential improvement, I believe accepting the opponent's proposal is the best course of action.} \\
\texttt{\}} \\
\\ 
\texttt{ARGUMENT: \{'Your proposal covers a wide range of items while efficiently using our available effort. The selected items provide a good balance across different categories and the effort utilization is well within our shared limit. This appears to be a well{-}thought{-}out selection that maximizes value while respecting our constraints.'\}} \\
\\ 
\texttt{AGREE: \{'C34', 'A27', 'C70', 'B45', 'B10', 'C35', 'A84', 'A30', 'B60', 'A89', 'A58', 'B09', 'A66', 'B46', 'A16', 'C92', 'A08', 'B94', 'B67', 'C81', 'B12', 'B00', 'B75'\}} \\
            }
        }
    }
     \\ \\

    \theutterance \stepcounter{utterance}  
    & & & \multicolumn{2}{p{0.3\linewidth}}{
        \cellcolor[rgb]{0.9,0.9,0.9}{
            \makecell[{{p{\linewidth}}}]{
                \texttt{\tiny{[GM$|$GM]}}
                \texttt{game successful} \\
            }
        }
    }
    & & \\ \\

    \theutterance \stepcounter{utterance}  
    & & & \multicolumn{2}{p{0.3\linewidth}}{
        \cellcolor[rgb]{0.9,0.9,0.9}{
            \makecell[{{p{\linewidth}}}]{
                \texttt{\tiny{[GM$|$GM]}}
                \texttt{end game} \\
            }
        }
    }
    & & \\ \\

\end{supertabular}
}

\end{document}
